\documentclass[11pt,a4paper]{article}

\usepackage[utf8]{inputenc}
\usepackage[T1]{fontenc}
\usepackage[brazil]{babel}
\usepackage[margin=2.5cm]{geometry}
\usepackage{amsmath}
\usepackage{graphicx}
\usepackage{booktabs}
\usepackage{siunitx}
\usepackage{float}
\usepackage{listings}
\usepackage{xcolor}
\usepackage[hidelinks]{hyperref}

\sisetup{output-decimal-marker={,}}
\graphicspath{{plots/}}

\lstset{
  basicstyle=\ttfamily\scriptsize,
  keywordstyle=\color{blue!60!black},
  commentstyle=\color{green!40!black},
  stringstyle=\color{red!60!black},
  numbers=left,
  numberstyle=\tiny\color{gray},
  breaklines=true,
  frame=single,
  language=Matlab,
  showstringspaces=false,
  captionpos=b
}

\title{Aplicações de Processamento Digital de Sinais -- 4456S-04\\[2mm]
\large Prática 6: Filtragem de Imagem}
\author{Antônio Vinicius Corrêa dos Santos}
\date{28/09/2026}

\begin{document}

\maketitle

Imagem: \texttt{lena\_c.jpg}, $256\times256$, convertida para escala de cinza com
\texttt{rgb2gray} e normalizada para $[0,1]$. Todas as figuras são geradas por
\texttt{main.m} e gravadas em \texttt{plots/}.

\section{Espectro de Fourier da imagem}

\begin{figure}[H]
\centering
\includegraphics[width=0.45\textwidth]{spectrum_original.jpg}
\caption{Espectro de magnitude, $\log(1+|\mathcal{F}\{x\}|)$, centralizado com
\texttt{fftshift}.}
\end{figure}

A energia se concentra no centro (baixas frequências) e cai rapidamente com o raio. As
duas faixas claras horizontal e vertical vêm da descontinuidade entre as bordas opostas
da imagem, que a DFT trata como periódica.

\section{Parte I --- filtragem no domínio espacial}

Cada máscara foi aplicada com \texttt{filter2} e o resultado exibido ao lado da imagem
original. A máscara passa-baixas foi normalizada por $1/9$ para manter o ganho DC
unitário. As saídas dos filtros derivativos são exibidas com \texttt{imshow(Y,[])},
porque têm valores negativos.

\begin{table}[H]
\centering
\caption{Máscaras utilizadas e efeito observado.}
\small
\begin{tabular}{llp{7.4cm}}
\toprule
\textbf{Máscara} & \textbf{Tipo} & \textbf{Efeito observado} \\
\midrule
Passa-baixas  & média $3\times3$ & Suavização leve; bordas e textura do cabelo perdem
definição. \\
Passa-altas   & soma nula, centro $8$ & Só as bordas sobrevivem; as regiões de
intensidade constante vão a zero. \\
Prewitt vert. / Sobel vert.  & derivada em $x$ & Realçam bordas verticais; Sobel dá
contornos mais fortes por causa do peso $2$ na linha central. \\
Prewitt hor. / Sobel hor.    & derivada em $y$ & O mesmo para bordas horizontais (chapéu,
ombro). \\
Kirsch        & direcional & Realce direcional de borda, com aspecto de relevo
(\emph{embossing}). \\
Laplaciano    & segunda derivada & Detecta borda nas duas direções, mas a resposta é
bem mais fraca que a dos de primeira derivada: só as transições mais abruptas aparecem,
como linhas finas sobre fundo uniforme. \\
\bottomrule
\end{tabular}
\end{table}

\begin{figure}[H]
\centering
\includegraphics[width=0.49\textwidth]{lena_c_Lowpass.jpg}
\hfill
\includegraphics[width=0.49\textwidth]{lena_c_Highpass.jpg}\\[1mm]
\includegraphics[width=0.49\textwidth]{lena_c_Sobel_Vertical.jpg}
\hfill
\includegraphics[width=0.49\textwidth]{lena_c_Sobel_Horizontal.jpg}\\[1mm]
\includegraphics[width=0.49\textwidth]{lena_c_Kirsch.jpg}
\hfill
\includegraphics[width=0.49\textwidth]{lena_c_Laplacian.jpg}
\caption{Original $\times$ filtrada: passa-baixas, passa-altas, Sobel vertical, Sobel
horizontal, Kirsch e Laplaciano. As saídas de Prewitt estão em
\texttt{plots/lena\_c\_Prewitt\_*.jpg} e são visualmente próximas das de Sobel.}
\end{figure}

\section{Parte II --- filtragem no domínio da frequência}

As máscaras são construídas sobre a matriz de distância $D$ centrada, então o produto é
feito com o espectro deslocado e desfeito antes da inversa:
\[
y = \Re\left\{ \mathcal{F}^{-1}\!\left[\text{ifftshift}\big(\text{fftshift}(\mathcal{F}\{x\})\cdot H\big)\right] \right\}.
\]
Multiplicar $H$ centrada diretamente por \texttt{fft2(img)} produz uma imagem
embaralhada; este foi o único ponto que exigiu atenção na implementação.

\[
H_{\text{ideal}}(u,v) = \begin{cases} 1, & D \le D_0 \\ 0, & D > D_0 \end{cases}
\qquad
H_{\text{butt}}(u,v) = \frac{1}{1 + (D/D_0)^{2n}}
\qquad
H_{\text{PA}} = 1 - H_{\text{butt}}
\]

\begin{figure}[H]
\centering
\includegraphics[width=0.85\textwidth]{masks_D0-50_n-2.jpg}
\caption{Máscaras para $D_0=50$, $n=2$.}
\end{figure}

\begin{figure}[H]
\centering
\includegraphics[width=0.49\textwidth]{lena_D0-50_n-2.jpg}
\hfill
\includegraphics[width=0.49\textwidth]{lena_D0-10_n-2.jpg}\\[1mm]
\includegraphics[width=0.49\textwidth]{lena_D0-120_n-2.jpg}
\hfill
\includegraphics[width=0.49\textwidth]{lena_D0-50_n-8.jpg}
\caption{Filtragem para $(D_0,n) = (50,2)$, $(10,2)$, $(120,2)$ e $(50,8)$. Em cada
figura: original, passa-baixas ideal, passa-baixas de Butterworth e passa-altas de
Butterworth.}
\end{figure}

\begin{table}[H]
\centering
\caption{Fração da energia do espectro que atravessa cada máscara,
$\sum |X\cdot H|^2 / \sum |X|^2$.}
\begin{tabular}{ccccc}
\toprule
$D_0$ & $n$ & \textbf{PB ideal} & \textbf{PB Butterworth} & \textbf{PA Butterworth} \\
\midrule
$50$  & $2$ & $0{,}9942$ & $0{,}9910$ & $0{,}0043$ \\
$10$  & $2$ & $0{,}9650$ & $0{,}9540$ & $0{,}0295$ \\
$120$ & $2$ & $0{,}9996$ & $0{,}9980$ & $0{,}0004$ \\
$50$  & $8$ & $0{,}9942$ & $0{,}9936$ & $0{,}0053$ \\
\bottomrule
\end{tabular}
\end{table}

Observações:

\begin{itemize}\itemsep0pt
  \item \textbf{Ideal $\times$ Butterworth.} O corte abrupto do filtro ideal produz
        \emph{ringing}: ondulações concêntricas em torno das bordas, bem visíveis em
        $D_0=10$. É o efeito de Gibbs em duas dimensões --- o equivalente espacial de
        $H_{\text{ideal}}$ é uma função tipo \emph{sinc} de suporte infinito. O
        Butterworth, com transição suave, não apresenta esse artefato; sua saída é
        apenas borrada.
  \item \textbf{Frequência de corte.} $D_0=10$ borra fortemente (resta pouco mais que a
        silhueta), $D_0=50$ suaviza a textura mantendo a forma, e $D_0=120$ quase não
        altera a imagem --- com $256\times256$, o raio de Nyquist é $128$, então
        $D_0=120$ deixa passar praticamente todo o espectro.
  \item \textbf{Ordem.} Aumentar $n$ de $2$ para $8$ com $D_0$ fixo aproxima o
        Butterworth do filtro ideal: a energia retida sobe de $0{,}9910$ para $0{,}9936$
        e começa a aparecer um leve \emph{ringing}, ainda muito menor que o do corte
        abrupto.
  \item \textbf{Passa-altas.} $1 - H_{\text{butt}}$ retém menos de \SI{3}{\percent} da
        energia e devolve apenas o mapa de bordas, coerente com o resultado das máscaras
        espaciais da Parte I.
\end{itemize}

\section{Imagem com \emph{blur} horizontal}

O \emph{blur} foi gerado com a máscara $h = \frac{1}{16}[1\;1\;\cdots\;1]$ ($1\times16$).

\begin{figure}[H]
\centering
\includegraphics[width=0.49\textwidth]{lena_c_Blurr_Horizontal.jpg}
\hfill
\includegraphics[width=0.36\textwidth]{spectrum_blur.jpg}
\caption{Imagem borrada na direção horizontal e seu espectro. As faixas escuras
verticais são os zeros do \emph{sinc} da média móvel.}
\end{figure}

O espectro mostra o que o \emph{blur} faz: a média móvel de 16 pontos tem resposta
$\sin(8\omega)/(16\sin(\omega/2))$, cujos zeros anulam completamente linhas inteiras de
frequências horizontais. Essa informação é perdida, não apenas atenuada.

\begin{table}[H]
\centering
\caption{Energia retida ao filtrar a imagem já borrada.}
\begin{tabular}{ccccc}
\toprule
$D_0$ & $n$ & \textbf{PB ideal} & \textbf{PB Butterworth} & \textbf{PA Butterworth} \\
\midrule
$50$  & $2$ & $0{,}9997$ & $0{,}9994$ & $0{,}0002$ \\
$10$  & $2$ & $0{,}9946$ & $0{,}9860$ & $0{,}0045$ \\
$120$ & $2$ & $1{,}0000$ & $0{,}9999$ & $0{,}0000$ \\
$50$  & $8$ & $0{,}9997$ & $0{,}9997$ & $0{,}0002$ \\
\bottomrule
\end{tabular}
\end{table}

Comparando com a tabela anterior, as frações são mais altas para os mesmos $(D_0,n)$:
a imagem borrada já não tem energia em alta frequência, então os passa-baixas quase não
têm o que remover. Aplicar passa-baixas sobre a imagem borrada só a borra mais;
nenhum deles reduz o \emph{blur}.

\subsection*{Tentativa de atenuar o \emph{blur}}

Como o \emph{blur} é uma perda de altas frequências, testamos realce por
\emph{unsharp masking}, $y = x_{\text{blur}} + k\,(1-H_{\text{butt}})\ast x_{\text{blur}}$,
variando $D_0$ e o ganho $k$.

\begin{figure}[H]
\centering
\includegraphics[width=0.85\textwidth]{blur_unsharp.jpg}
\caption{Imagem borrada e duas tentativas de realce.}
\end{figure}

\begin{table}[H]
\centering
\caption{Erro RMS em relação à imagem original.}
\begin{tabular}{lc}
\toprule
\textbf{Imagem} & \textbf{RMS} \\
\midrule
Borrada                          & $0{,}0944$ \\
Realçada, $D_0=10$, $n=2$, $k=0{,}25$ & $0{,}0939$ \\
Realçada, $D_0=5$,  $n=2$, $k=1{,}0$  & $0{,}1116$ \\
\bottomrule
\end{tabular}
\end{table}

O realce aumenta o contraste das bordas que sobreviveram e a imagem parece mais nítida,
mas o erro em relação à original praticamente não melhora ($0{,}0944 \to 0{,}0939$) e
piora com ganho alto. A razão é que as máscaras são \emph{isotrópicas} e o \emph{blur} é
direcional: o realce amplifica também as frequências verticais, que nunca foram
degradadas, enquanto as horizontais anuladas pelos zeros do \emph{sinc} não podem ser
recuperadas por multiplicação no domínio da frequência --- seriam necessárias divisão por
$H$ (filtro inverso) ou um filtro de Wiener, fora do escopo desta prática. A conclusão da
parte final é negativa e esperada: filtragem passa-baixas/passa-altas isotrópica não
desfaz um \emph{blur} direcional.

\section*{Script (Octave)}

\lstinputlisting[caption={\texttt{main.m}}]{main.m}

\end{document}
